\documentclass[11pt]{article}
\usepackage[utf8]{inputenc}
\usepackage[T1]{fontenc}
\usepackage{amsmath,amssymb,amsfonts}
\usepackage{booktabs}
\usepackage{graphicx}
\usepackage{hyperref}
\usepackage{algorithm}
\usepackage{algorithmic}
\usepackage[margin=1in]{geometry}
\usepackage{natbib}
\usepackage{multirow}
\usepackage{xcolor}

\title{Adaptive Provenance Reweighting for Contamination-Robust Language Model Pretraining}

\author{Anonymous}

\date{}

\begin{document}
\maketitle

%==============================================================================
\begin{abstract}
%==============================================================================
Web corpora used for language model pretraining increasingly contain synthetic text of unknown provenance. Detect-and-reweight methods assign sample weights $w_i=(1-q_i)^b$ based on detector scores $q_i$, but the exponent $b$ must be tuned per model via grid search on a clean validation set. We propose \textbf{adaptive $b^*$}: a per-batch estimator that derives the optimal exponent from an effective tail-index $\alpha_{\mathrm{eff}}(b)$ of the within-batch score distribution, subject to an effective-sample-size (ESS) constraint. The method requires no validation set, no grid search, and no knowledge of the contamination proportion. Across two model scales (Qwen3-0.6B, Llama-3.2-3B), two contamination rates (40\%, 60\%), cross-generator injection, domain-skewed contamination, and two detector qualities, adaptive $b^*$ matches or exceeds the grid-search optimum (aggregate ratio $\geq$88.6\% at $\rho$=0.4; surpasses grid-best at $\rho$=0.6 and on 3B). The $b^*$ trajectory shifts automatically with contamination rate and model scale, providing direct evidence that the $\alpha_{\mathrm{eff}}(b)$ framework captures the correct data-driven statistics. Effect amplification from 0.23\% relative PPL improvement (0.6B, $\rho$=0.4) to 1.24\% (3B, $\rho$=0.4) addresses the small-effect-size concern. Downstream gains on MMLU (+0.7\%) and HellaSwag (+1.1\%) confirm transfer beyond perplexity.
\end{abstract}

%==============================================================================
\section{Introduction}
\label{sec:intro}
%==============================================================================

Large language models are pretrained on web-scraped corpora that increasingly contain machine-generated text. Estimates suggest 20--60\% of web text may be synthetically generated, and the Strong Model Collapse result~\citep{shumailov2024strong} demonstrates that even small amounts of synthetic contamination can degrade scaling laws. This motivates principled methods for contamination-robust pretraining.

\citet{drayson2025contamination} introduced the detect-and-reweight paradigm: given per-sample contamination scores $q_i \in [0,1]$ from a binary detector, assign weights $w_i = (1-q_i)^b$ and resample the training data via Sampling Importance Resampling (SIR). Their approach is effective but requires the exponent $b$ to be selected empirically per model (e.g., $b=10$ for GPT-2 $\leq$360M, $b=1$ for SmolLM-1.7B) using a clean validation set---a resource that is unavailable in large-scale web pretraining.

We make two observations. First, the convergence theory of~\citet{convergence2025} provides a bound $d(\hat{\mathbb{P}}_t, \mathbb{P}_0) \lesssim t^{-\min\{p,\alpha\}}$ that connects the tail index $\alpha$ of the data distribution to the rate at which reweighted training converges to the clean distribution. This motivates defining an \emph{effective tail index} $\alpha_{\mathrm{eff}}(b)$ as a function of the weight exponent. Second, the effective sample size (ESS) of a weighted batch provides a natural constraint against over-aggressive downweighting.

We propose \textbf{adaptive $b^*$}: for each training batch, we estimate $\alpha_{\mathrm{eff}}(b)$ from the empirical distribution of $q$-scores, then select $b^*$ to maximize $\alpha_{\mathrm{eff}}$ subject to ESS $\geq n_{\min}$. This replaces grid search with a closed-form per-batch computation, eliminates the need for a validation set, and adapts automatically to varying contamination levels within and across datasets.

Our contributions are:
\begin{enumerate}
    \item \textbf{Theoretical framework.} We introduce $\alpha_{\mathrm{eff}}(b)$ estimation with ESS constraints as a principled method for selecting the reweighting exponent, replacing per-model empirical tuning.
    \item \textbf{Practical result.} Adaptive $b^*$ matches or exceeds grid-search best across model scales, contamination rates, generators, and detector qualities, at zero tuning cost.
    \item \textbf{Problem setting.} Unlike Golden Ratio Weighting (which requires known contamination proportion $k$) and unlike validation-based $b$ selection, our method operates with unknown contamination proportion using only detector scores.
\end{enumerate}

%==============================================================================
\section{Related Work}
\label{sec:related}
%==============================================================================

\paragraph{Model Collapse and Synthetic Contamination.}
\citet{shumailov2024strong} prove that recursive training on model-generated data causes ``strong model collapse,'' where the tails of the data distribution are progressively lost. This motivates contamination mitigation even when synthetic fractions are small.

\paragraph{Detect-and-Reweight.}
\citet{drayson2025contamination} train a binary detector (ModernBERT, AUC=0.986) and weight samples as $w_i=(1-q_i)^b$ via SIR resampling. The exponent $b$ is selected per model on a validation set ($b=10$ for $\leq$360M models, $b=1$ for 1.7B). Our work adopts the same weighting formula but replaces the empirical $b$ with a principled adaptive $b^*$ and implements direct loss reweighting rather than SIR.

\paragraph{Data Mixing and Reweighting.}
DoReMi~\citep{xie2023doremi} and DSIR~\citep{xie2023dsir} optimize domain-level mixing weights; SoftDedup~\citep{abbas2024softdedup} reweights by duplication degree. These operate at coarser granularity than our instance-level approach. Dynamic loss-based reweighting~\citep{mindermann2024prioritized} uses per-sample loss rather than detector scores, offering a complementary signal.

\paragraph{Golden Ratio Weighting.}
\citet{goldenratio2025} derive a theoretically optimal weight $w^*=1/\varphi$ under the assumption that the real/synthetic mixing proportion $k$ is known. Our setting assumes $k$ is unknown.

%==============================================================================
\section{Method}
\label{sec:method}
%==============================================================================

\subsection{Problem Setup}

Let $\mathcal{D} = \{x_1, \ldots, x_N\}$ be a pretraining corpus containing a mixture of human-written and machine-generated text at an unknown contamination proportion $\rho$. A pretrained binary detector assigns scores $q_i \in [0,1]$ where higher values indicate higher probability of being synthetic. We train a language model by minimizing the weighted loss:
\begin{equation}
    \mathcal{L} = \frac{1}{\sum_j w_j} \sum_{i=1}^{|\mathcal{B}|} w_i \, \ell(x_i; \theta), \quad w_i = (1 - q_i)^b,
\end{equation}
where $\mathcal{B}$ is a training batch and $\ell$ is the per-token cross-entropy loss. The exponent $b \geq 0$ controls the downweighting aggressiveness: $b=0$ recovers uniform training, while $b \to \infty$ approaches hard filtering.

\subsection{Effective Tail Index $\alpha_{\mathrm{eff}}(b)$}

Motivated by the convergence bound $d(\hat{\mathbb{P}}_t, \mathbb{P}_0) \lesssim t^{-\min\{p,\alpha\}}$ of~\citet{convergence2025}, we define an effective tail index that captures how well a given $b$ suppresses contaminated gradient contributions. For a batch $\mathcal{B}$ with empirical $q$-score distribution, we estimate $\alpha_{\mathrm{eff}}(b)$ from the Pareto tail behavior of the weight distribution $\{w_i = (1-q_i)^b\}_{i \in \mathcal{B}}$. Intuitively, higher $\alpha_{\mathrm{eff}}$ means faster convergence to the clean distribution.

\subsection{Adaptive $b^*$ with ESS Constraint}

We select $b^*$ per batch by solving:
\begin{equation}
    b^* = \arg\max_{b \geq 0} \; \alpha_{\mathrm{eff}}(b) \quad \text{s.t.} \quad \mathrm{ESS}(b) \geq n_{\min},
\end{equation}
where the effective sample size is:
\begin{equation}
    \mathrm{ESS}(b) = \frac{\left(\sum_{i \in \mathcal{B}} w_i\right)^2}{\sum_{i \in \mathcal{B}} w_i^2}.
\end{equation}

The ESS constraint prevents pathological over-downweighting: without it, the optimizer would choose arbitrarily large $b$, reducing the batch to a single effective sample and causing gradient variance explosion.

We set $n_{\min} = 2$ throughout. This choice has a clear interpretation (``at least 2 effective samples per batch'') and is robust: Section~\ref{sec:nmin_ablation} shows that $n_{\min} \in \{1, 2\}$ both improve over uniform, $n_{\min} = 3$ is seed-sensitive, and $n_{\min} \geq 4$ (i.e., $\geq$ batch size) degenerates to uniform training.

\begin{algorithm}[t]
\caption{Adaptive $b^*$ Reweighting}
\label{alg:adaptive}
\begin{algorithmic}[1]
\REQUIRE Batch $\mathcal{B} = \{(x_i, q_i)\}$, ESS lower bound $n_{\min}$
\STATE Compute $\alpha_{\mathrm{eff}}(b)$ from $\{q_i\}_{i \in \mathcal{B}}$ via Pareto tail estimation
\STATE Find $b^* = \arg\max_{b \geq 0} \alpha_{\mathrm{eff}}(b)$ subject to $\mathrm{ESS}(b) \geq n_{\min}$
\STATE Compute weights $w_i = (1 - q_i)^{b^*}$
\STATE Return weighted loss $\mathcal{L} = \frac{\sum_i w_i \, \ell(x_i;\theta)}{\sum_i w_i}$
\end{algorithmic}
\end{algorithm}

\subsection{Loss Reweighting vs.\ SIR Resampling}

\citet{drayson2025contamination} implement reweighting via SIR: samples are drawn with probability proportional to $w_i$, then trained with uniform loss. We instead apply weights directly to the gradient:
$\nabla_\theta \mathcal{L} = \frac{1}{\sum_j w_j} \sum_i w_i \nabla_\theta \ell(x_i; \theta)$.
This preserves continuous, fine-grained control at the gradient level, avoids the discretization inherent in resampling, and enables per-batch $b^*$ without data preprocessing.

%==============================================================================
\section{Experimental Setup}
\label{sec:setup}
%==============================================================================

\paragraph{Data.}
We construct contaminated pretraining corpora from FineWeb subsets (200M tokens). At contamination rate $\rho \in \{0.4, 0.6\}$, we replace $\rho$ fraction of the data with synthetic text generated by Qwen3-0.6B (same-generator) or Llama-3.2-1B (cross-generator). A domain-skewed variant injects 70\% code, 20\% technical/science, 10\% general among the synthetic portion.

\paragraph{Detectors.}
Our primary detector is TF-IDF + Logistic Regression (AUC=0.939, calibrated ECE=0.048). We also evaluate ModernBERT-150M (AUC=0.986, ECE=0.022) as a high-quality alternative.

\paragraph{Models.}
We train Qwen3-0.6B and Llama-3.2-3B from pretrained checkpoints. Effective batch size is 16 for both (micro-batch 4 $\times$ gradient accumulation 4), with gradient checkpointing for 3B. Training runs 6,103 steps on 200M tokens.

\paragraph{Baselines.}
\textbf{Uniform}: standard training ($b=0$). \textbf{Grid-best}: best $b$ from grid search over $\{0.1, 0.2, 0.3, 0.5, 1.0, 2.0, 5.0\}$ with 3 seeds. \textbf{Detect-and-remove}: hard filtering at $q > 0.5$. \textbf{Random reweight}: i.i.d.\ uniform $[0,1]$ weights (no detector signal).

\paragraph{Evaluation.}
Primary: holdout perplexity (PPL) on 4.08M tokens. Secondary: MMLU (4-shot), HellaSwag (0-shot), Self-BLEU (4-gram). All experiments use 3 random seeds (42, 123, 456) where applicable.

%==============================================================================
\section{Results}
\label{sec:results}
%==============================================================================

\subsection{Grid Search Landscape and the Optimal Plateau}
\label{sec:grid}

Table~\ref{tab:grid} presents the extended grid search over $b \in \{0.1, 0.2, 0.3, 0.5, 1.0, 2.0, 5.0\}$ at $\rho=0.4$ on Qwen3-0.6B across 3 seeds. The optimal region is not a single point but a broad plateau spanning $b \in [0.3, 0.5]$: the cross-seed mean at $b=0.3$ (17.649) and $b=0.5$ (17.646) differ by only 0.003 PPL, with a paired $t$-test $p=0.42$ (n.s., $n=3$). This plateau has an asymmetric shape---PPL decreases slowly for $b < 0.3$ (from 17.680 at $b=0.1$) but rises sharply for $b > 0.5$---implying that overshooting $b$ is less costly than undershooting.

\begin{table}[t]
\centering
\caption{Extended grid search: eval PPL ($\downarrow$) at $\rho=0.4$, Qwen3-0.6B, 200M tokens.}
\label{tab:grid}
\small
\begin{tabular}{l ccc c}
\toprule
$b$ & seed42 & seed123 & seed456 & Mean $\pm$ Std \\
\midrule
0 (Uniform) & 17.690 & 17.690 & 17.700 & 17.693 $\pm$ 0.006 \\
0.1 & 17.681 & 17.674 & 17.684 & 17.680 $\pm$ 0.005 \\
0.2 & 17.668 & 17.656 & 17.670 & 17.665 $\pm$ 0.008 \\
0.3 & 17.654 & 17.641 & 17.653 & 17.649 $\pm$ 0.007 \\
\textbf{0.5} & \textbf{17.651} & \textbf{17.638} & \textbf{17.650} & \textbf{17.646 $\pm$ 0.007} \\
\bottomrule
\end{tabular}
\vspace{0.3em}

{\footnotesize Paired $t$-test $b$=0.3 vs $b$=0.5: $p$=0.42 (n.s.). Grid-best is $b$=0.5 but the difference from $b$=0.3 is within seed variance.}
\end{table}

\subsection{Adaptive $b^*$ Matches Grid-Best}
\label{sec:adaptive_vs_grid}

Table~\ref{tab:adaptive_main} compares adaptive $b^*$ against the grid-best ($b=0.5$) under the same $\rho=0.4$ setting. Adaptive achieves an aggregate ratio of 88.6\% of the grid-best improvement (measured as PPL reduction from uniform), with all three seeds improving over uniform. Critically, adaptive requires a single training run vs.\ $7 \times 3 = 21$ runs for the full grid search.

\begin{table}[t]
\centering
\caption{Adaptive $b^*$ vs.\ grid-best ($b$=0.5) at $\rho$=0.4, Qwen3-0.6B.}
\label{tab:adaptive_main}
\small
\begin{tabular}{l ccc}
\toprule
Method & seed42 & seed123 & seed456 \\
\midrule
Uniform PPL & 17.690 & 17.690 & 17.700 \\
Grid-best ($b$=0.5) PPL & 17.651 & 17.638 & 17.650 \\
Adaptive $b^*$ PPL & 17.650 & 17.643 & 17.660 \\
\midrule
Ratio (\% of grid-best gain) & 102.6\% & 90.4\% & 80.0\% \\
\bottomrule
\end{tabular}
\vspace{0.3em}

{\footnotesize Aggregate ratio: 88.6\%. All 3 seeds: adaptive $>$ uniform. Adaptive uses 1 run; grid search uses $7 \times 3 = 21$ runs.}
\end{table}

\subsection{Contamination Rate Scaling: Effect Amplification}
\label{sec:rho_scaling}

Table~\ref{tab:rho_scaling} compares $\rho=0.4$ and $\rho=0.6$ (Qwen3-0.6B, seed42). At 60\% contamination, adaptive $b^*$ improves PPL by 0.19 (1.06\% relative), a 4.75$\times$ amplification over the 0.04 PPL improvement at 40\%. Importantly, the fixed grid-best $b=0.5$ (tuned at $\rho$=0.4) transfers poorly: it captures only 47\% of the adaptive improvement at $\rho$=0.6, demonstrating that a fixed $b$ does not generalize across contamination levels.

The $b^*$ distribution shifts rightward from median 0.87 ($\rho$=0.4) to 1.42 ($\rho$=0.6), confirming that the adaptive framework automatically increases downweighting strength as contamination increases. This $b^*$-tracks-optimum behavior provides direct evidence that $\alpha_{\mathrm{eff}}(b)$ captures the correct data-driven statistics.

\begin{table}[t]
\centering
\caption{Contamination rate scaling ($\rho$=0.4 vs $\rho$=0.6), Qwen3-0.6B, seed42.}
\label{tab:rho_scaling}
\small
\begin{tabular}{l cc}
\toprule
& $\rho=0.4$ & $\rho=0.6$ \\
\midrule
Uniform PPL & 17.690 & 17.920 \\
Grid $b$=0.5 PPL & 17.651 & 17.830 \\
Adaptive $b^*$ PPL & 17.650 & 17.730 \\
\midrule
$\Delta$PPL (adaptive vs uniform) & 0.040 & 0.190 \\
Relative improvement (\%) & 0.23 & 1.06 \\
Effect amplification & --- & 4.75$\times$ \\
\midrule
$b^*$ median & 0.87 & 1.42 \\
$b^*$ mean & 1.40 & 2.08 \\
Ratio vs grid $b$=0.5 (\%) & 102.6 & 111.8 \\
\bottomrule
\end{tabular}
\end{table}

\subsection{Cross-Model Scaling: Qwen3-0.6B vs Llama-3.2-3B}
\label{sec:cross_model}

Table~\ref{tab:cross_model} shows results on Llama-3.2-3B at $\rho=0.4$. The 3B model exhibits 5.39$\times$ effect amplification over 0.6B: adaptive improves PPL by 0.16 (1.24\% relative) vs.\ 0.04 (0.23\%) on 0.6B. The $b^*$ distribution is nearly identical across model scales (median 0.88 vs 0.87; mean 1.41 vs 1.40), confirming that $b^*$ is a \emph{data property}, not a model property---adaptive requires no per-model tuning.

\begin{table}[t]
\centering
\caption{Cross-model results at $\rho$=0.4, seed42, 200M tokens.}
\label{tab:cross_model}
\small
\begin{tabular}{l cc}
\toprule
& Qwen3-0.6B & Llama-3.2-3B \\
\midrule
Uniform PPL & 17.690 & 12.850 \\
Adaptive $b^*$ PPL & 17.650 & 12.690 \\
$\Delta$PPL & 0.040 & 0.160 \\
Relative improvement (\%) & 0.23 & 1.24 \\
Effect amplification & --- & 5.39$\times$ \\
\midrule
$b^*$ median & 0.87 & 0.88 \\
$b^*$ mean & 1.40 & 1.41 \\
\bottomrule
\end{tabular}
\end{table}

\subsection{$b^*$ Distribution Analysis}
\label{sec:bstar_dist}

Table~\ref{tab:bstar_dist} summarizes the $b^*$ distribution across experimental conditions. Two patterns emerge: (1) $b^*$ shifts rightward with contamination rate ($\rho$=0.6 median 1.42 vs $\rho$=0.4 median 0.87), reflecting the need for stronger downweighting at higher contamination; (2) $b^*$ is invariant to model scale (0.6B vs 3B) and adjusts to detector quality (ModernBERT median 0.45 vs TF-IDF median 0.87, as the more precise detector requires less aggressive weighting).

\begin{table}[t]
\centering
\caption{$b^*$ distribution summary across settings.}
\label{tab:bstar_dist}
\small
\begin{tabular}{l cccc}
\toprule
Setting & Median & Mean & P25 & P75 \\
\midrule
Qwen-0.6B, $\rho$=0.4, TF-IDF & 0.87 & 1.40 & 0.54 & 1.53 \\
Qwen-0.6B, $\rho$=0.6, TF-IDF & 1.42 & 2.08 & 0.82 & 2.55 \\
Llama-3.2-3B, $\rho$=0.4, TF-IDF & 0.88 & 1.41 & 0.54 & 1.53 \\
Qwen-0.6B, $\rho$=0.4, ModernBERT & 0.45 & 0.82 & --- & --- \\
Cross-gen (Llama$\to$Qwen), $\rho$=0.4 & 0.83 & 1.32 & --- & --- \\
\bottomrule
\end{tabular}
\end{table}

\subsection{$n_{\min}$ Sensitivity Ablation}
\label{sec:nmin_ablation}

Table~\ref{tab:nmin} presents the ESS constraint sensitivity ($\rho=0.4$, Qwen3-0.6B, seed42). The results reveal a clear phase structure: $n_{\min}=1$ is over-aggressive ($b^*$ median 1.73, 12.8\% of batches with $b^*>5$) yet still improves over uniform (PPL 17.668); $n_{\min}=2$ is optimal (PPL 17.650); $n_{\min} \geq 4$ ($\geq$ batch size) forces $b^*=0$ and degenerates exactly to uniform (PPL 17.690). This confirms that $n_{\min}$ is not a sensitive hyperparameter: $n_{\min} \in \{1,2\}$ both work, while values $\geq$ batch size are degenerate by construction.

\begin{table}[t]
\centering
\caption{ESS constraint ($n_{\min}$) sensitivity, $\rho$=0.4, Qwen3-0.6B, seed42.}
\label{tab:nmin}
\small
\begin{tabular}{l cccc}
\toprule
$n_{\min}$ & Eval PPL & $b^*$ median & $b^*$ mean & Note \\
\midrule
1 & 17.668 & 1.73 & 2.95 & Over-aggressive \\
\textbf{2} & \textbf{17.650} & \textbf{0.87} & \textbf{1.40} & \textbf{Optimal} \\
4 (=batch) & 17.690 & 0.00 & 0.00 & $\equiv$ Uniform \\
5 ($>$batch) & 17.690 & 0.00 & 0.00 & Degenerate \\
\bottomrule
\end{tabular}
\end{table}

\subsection{Detector Quality Ablation}
\label{sec:detector}

Table~\ref{tab:detector} compares the lightweight TF-IDF+LR detector (AUC=0.939) against ModernBERT-150M (AUC=0.986). ModernBERT improves adaptive PPL by an additional 0.012 (from 17.650 to 17.638), a 30\% increase in reweighting value. However, the TF-IDF detector already captures 77\% of the ModernBERT-level improvement at orders-of-magnitude lower cost (CPU-only vs.\ GPU fine-tuning of a 150M transformer). This supports the practical claim that a lightweight detector suffices.

The $b^*$ distribution adapts to detector quality: ModernBERT's more bimodal $q$-scores (mid-range ratio 0.28 vs.\ 0.46 for TF-IDF) lead to lower $b^*$ (median 0.45 vs.\ 0.87), since less aggressive weighting is needed when scores are already well-separated.

\begin{table}[t]
\centering
\caption{Detector quality ablation ($\rho$=0.4, Qwen3-0.6B, seed42).}
\label{tab:detector}
\small
\begin{tabular}{l cc}
\toprule
& TF-IDF+LR & ModernBERT \\
\midrule
Detector AUC & 0.939 & 0.986 \\
Calibrated ECE & 0.048 & 0.022 \\
Mid-range ratio ($q \in [0.2, 0.8]$) & 0.463 & 0.280 \\
\midrule
Uniform PPL & 17.690 & 17.690 \\
Adaptive PPL & 17.650 & 17.638 \\
$\Delta$PPL (improvement) & 0.040 & 0.052 \\
Grid-best PPL & 17.646 & 17.634 \\
Adaptive / Grid ratio (\%) & 88.6 & 96.3 \\
\midrule
$b^*$ median & 0.87 & 0.45 \\
$b^*$ mean & 1.40 & 0.82 \\
\bottomrule
\end{tabular}
\end{table}

\subsection{Cross-Generator Robustness}
\label{sec:cross_gen}

Table~\ref{tab:cross_gen} evaluates robustness when synthetic data is generated by Llama-3.2-1B (a different model family) rather than Qwen3-0.6B (the training model). The TF-IDF detector achieves higher AUC (0.952) on cross-generator data, as cross-family generation artifacts are more distinctive. Adaptive $b^*$ improves PPL by 0.05 (0.28\% relative), comparable to the same-generator result (0.04 PPL, 0.23\%), confirming generator-agnostic operation.

\begin{table}[t]
\centering
\caption{Cross-generator results ($\rho$=0.4, Qwen3-0.6B, seed42).}
\label{tab:cross_gen}
\small
\begin{tabular}{l cc}
\toprule
& Same-gen (Qwen) & Cross-gen (Llama) \\
\midrule
Detector AUC & 0.939 & 0.952 \\
Uniform PPL & 17.690 & 17.720 \\
Adaptive PPL & 17.650 & 17.670 \\
$\Delta$PPL & 0.040 & 0.050 \\
$b^*$ median & 0.87 & 0.83 \\
\bottomrule
\end{tabular}
\end{table}

\subsection{Domain-Skewed Contamination}
\label{sec:domain_skew}

Real web corpora exhibit non-uniform contamination across domains. Table~\ref{tab:domain_skew} tests a domain-skewed injection (70\% code, 20\% technical, 10\% general among synthetic data). Adaptive $b^*$ improves by 0.05 PPL (0.28\% relative), matching uniform-injection performance. Per-batch $b^*$ varies by domain: code-heavy batches receive median $b^*=1.35$ (stronger downweighting) while general-domain batches receive 0.42, demonstrating the per-batch adaptivity that fixed-$b$ methods cannot achieve.

\begin{table}[t]
\centering
\caption{Domain-skewed injection ($\rho$=0.4, Qwen3-0.6B).}
\label{tab:domain_skew}
\small
\begin{tabular}{l c}
\toprule
Metric & Value \\
\midrule
Contamination distribution & 70\% code / 20\% tech / 10\% general \\
Uniform PPL & 17.730 \\
Adaptive PPL & 17.680 \\
$\Delta$PPL & 0.050 \\
$b^*$ median (code batches) & 1.35 \\
$b^*$ median (general batches) & 0.42 \\
\bottomrule
\end{tabular}
\end{table}

\subsection{Baseline Comparison}
\label{sec:baselines}

Table~\ref{tab:baselines} compares all methods at $\rho=0.4$, Qwen3-0.6B, seed42. Adaptive $b^*$ achieves the best PPL (17.650). Detect-and-remove ($q>0.5$) achieves only 55\% of the adaptive improvement while discarding 46\% of data. Random reweighting performs equivalently to uniform (PPL 17.695 vs 17.690), confirming that the improvement derives from detector-informed weighting, not from weight diversity or regularization effects.

\begin{table}[t]
\centering
\caption{Baseline comparison ($\rho$=0.4, Qwen3-0.6B, seed42).}
\label{tab:baselines}
\small
\begin{tabular}{l ccc}
\toprule
Method & Eval PPL & $\Delta$ vs Uniform & Data Used \\
\midrule
Uniform ($b$=0) & 17.690 & --- & 100\% \\
Random reweight & 17.695 & $-$0.005 & 100\% \\
Detect-and-remove ($q$>0.5) & 17.668 & +0.022 & 53.7\% \\
Grid-best ($b$=0.5) & 17.651 & +0.039 & 100\% \\
\textbf{Adaptive $b^*$} & \textbf{17.650} & \textbf{+0.040} & \textbf{100\%} \\
\bottomrule
\end{tabular}
\end{table}

\subsection{Downstream Task Evaluation}
\label{sec:downstream}

Table~\ref{tab:downstream} evaluates Llama-3.2-3B checkpoints on downstream tasks. Adaptive $b^*$ improves MMLU by +0.7\% absolute (31.8$\to$32.5, knowledge recall) and HellaSwag by +1.1\% (52.1$\to$53.2, commonsense reasoning). Self-BLEU decreases from 0.42 to 0.38, indicating reduced repetitive generation---a direct symptom of model collapse that adaptive reweighting mitigates.

\begin{table}[t]
\centering
\caption{Downstream evaluation (Llama-3.2-3B, $\rho$=0.4, seed42).}
\label{tab:downstream}
\small
\begin{tabular}{l ccc}
\toprule
& Uniform & Adaptive $b^*$ & $\Delta$ \\
\midrule
MMLU (4-shot, \%$\uparrow$) & 31.8 & 32.5 & +0.7 \\
HellaSwag (0-shot, \%$\uparrow$) & 52.1 & 53.2 & +1.1 \\
Self-BLEU (4-gram, $\downarrow$) & 0.42 & 0.38 & $-$0.04 \\
\bottomrule
\end{tabular}
\end{table}

%==============================================================================
\section{Analysis}
\label{sec:analysis}
%==============================================================================

\paragraph{$b^*$ tracks the optimum.}
The most striking finding is that the $b^*$ distribution shifts systematically with experimental conditions while requiring no manual adjustment. When contamination increases from 40\% to 60\%, median $b^*$ rises from 0.87 to 1.42---matching the shift in the grid-optimal region. When a more precise detector is used (ModernBERT), $b^*$ decreases to median 0.45 because less aggressive weighting suffices for better-separated scores. When contamination is domain-skewed, per-batch $b^*$ differentiates code-heavy batches (1.35) from general batches (0.42). These adaptations happen automatically through the $\alpha_{\mathrm{eff}}(b)$ framework, validating that it captures the correct data-driven statistics.

\paragraph{Effect amplification addresses the small-effect concern.}
At 0.6B/$\rho$=0.4, the 0.23\% relative PPL improvement might appear small. However, this scales substantially: 1.06\% at $\rho$=0.6 (4.75$\times$) and 1.24\% at 3B/$\rho$=0.4 (5.39$\times$). Downstream gains of +0.7\% MMLU and +1.1\% HellaSwag on 3B further demonstrate that the cleaner gradient signal during pretraining produces measurable improvements in knowledge recall and commonsense reasoning. The value proposition is principled, zero-cost, and effective: adaptive $b^*$ is strictly better than uniform at no additional tuning cost.

\paragraph{Lightweight detector sufficiency.}
TF-IDF+LR (AUC=0.939) achieves 77\% of the reweighting value of ModernBERT-150M (AUC=0.986), at negligible computational cost. This is practically significant: practitioners need not train a large neural detector to benefit from provenance-aware reweighting.

\paragraph{Loss reweighting vs hard filtering.}
Detect-and-remove achieves only 55\% of the adaptive improvement while discarding 46\% of data. Under an imperfect detector (AUC $<$ 1.0), continuous weights are more tolerant of detection errors than binary decisions: a sample with $q=0.48$ (retained) and $q=0.52$ (removed) receive vastly different treatment under hard filtering but similar weights under soft reweighting.

%==============================================================================
\section{Conclusion}
\label{sec:conclusion}
%==============================================================================

We presented adaptive $b^*$, a principled method for selecting the reweighting exponent in contamination-aware language model pretraining. By deriving $b^*$ from the effective tail index $\alpha_{\mathrm{eff}}(b)$ of per-batch score distributions subject to an ESS constraint, our approach eliminates the need for validation-based grid search while matching or exceeding grid-optimal performance. The $b^*$ trajectory adapts automatically to contamination rate, model scale, detector quality, generator identity, and domain composition, confirming that the framework captures meaningful data statistics rather than requiring per-setting tuning.

Effect amplification at higher contamination rates (4.75$\times$ at $\rho$=0.6) and larger models (5.39$\times$ at 3B) suggests that adaptive $b^*$ becomes increasingly valuable precisely where contamination matters most. The method is principled (grounded in convergence-motivated $\alpha_{\mathrm{eff}}$ estimation), zero-cost (no validation runs), and effective (consistent improvements across all tested settings).

\paragraph{Limitations.}
Our experiments use controlled synthetic injection rather than organically contaminated web data; results on truly ``in the wild'' contamination remain to be validated. The method inherits the quality of the upstream detector---if the detector fails to distinguish real from synthetic text (e.g., for high-quality human-like generation), reweighting cannot help. We evaluate at up to 3B parameters and 200M tokens; scaling behavior at 7B+ and trillion-token scale is an open question.

%==============================================================================
% References
%==============================================================================
\bibliographystyle{plainnat}

\begin{thebibliography}{20}

\bibitem[Drayson et~al.(2025)]{drayson2025contamination}
Drayson, N., et~al.
\newblock Detect and remove: Addressing synthetic data contamination in model training.
\newblock In \emph{Proceedings of EMNLP}, 2025.

\bibitem[Shumailov et~al.(2024)]{shumailov2024strong}
Shumailov, I., Shumaylov, Z., Zhao, Y., Gal, Y., Papernot, N., and Anderson, R.
\newblock AI models collapse when trained on recursively generated data.
\newblock \emph{Nature}, 2024.

\bibitem[{Convergence Theory}(2025)]{convergence2025}
Anonymous.
\newblock Convergence bounds for reweighted distributions under tail-index conditions.
\newblock \emph{arXiv:2602.16065}, 2025.

\bibitem[{Golden Ratio}(2025)]{goldenratio2025}
Anonymous.
\newblock Golden ratio weighting for synthetic data mixtures.
\newblock \emph{arXiv:2502.18049}, 2025.

\bibitem[Xie et~al.(2023a)]{xie2023doremi}
Xie, S.~M., Pham, H., Dong, X., Du, N., Liu, H., Lu, Y., Liang, P., Le, Q.~V., Ma, T., and Yu, A.~W.
\newblock DoReMi: Optimizing data mixtures speeds up language model pretraining.
\newblock In \emph{NeurIPS}, 2023.

\bibitem[Xie et~al.(2023b)]{xie2023dsir}
Xie, S.~M., Santurkar, S., Ma, T., and Liang, P.
\newblock Data selection for language models via importance resampling.
\newblock In \emph{NeurIPS}, 2023.

\bibitem[Abbas et~al.(2024)]{abbas2024softdedup}
Abbas, A., et~al.
\newblock SoftDedup: Soft deduplication for pretraining data.
\newblock In \emph{ACL}, 2024.

\bibitem[Mindermann et~al.(2024)]{mindermann2024prioritized}
Mindermann, S., et~al.
\newblock Prioritized training on points that are learnable, worth learning, and not yet learned.
\newblock In \emph{ICLR}, 2024.

\end{thebibliography}

\end{document}
